\section{The deformation inequality for a single character}
\label{sec:single-character}

The estimate used in Section~\ref{sec:rh-plus-nonzero} first bounds one
strict Selmer group by distances between traces. We give the lower
comparison with the saturated cotangent module, the multiplicities in
the upper comparison, and the passage to isomorphism classes. We state
which deformation and automorphy results enter these comparisons.

\subsection{The arithmetic hypotheses and the estimate}

Take an algebraic closure \(\Omega/\mathbb Q\), an imaginary quadratic
subfield \(K\), a prime \(\ell\), and the integers \(\mathcal O\) in a
finite extension \(F/\mathbb Q_\ell\) inside \(\overline{\mathbb Q}_\ell\).
Write \(e_\mathcal O=v_{\mathcal O}(\ell)\), with
\(v_{\mathcal O}(\varpi)=1\). Fix an even positive integer \(N\), an
integer \(e\ge1\) and an \(\mathcal O\)-valued Dirichlet character
\(\psi\bmod N\). Write \(\psi_\ell\) for its component at the
\(\ell\)-part of the level and \(\psi'=\psi/\psi_\ell\) on the units.
The good-data assumptions of the source are
\[
\begin{gathered}
 \ell\ne2,\qquad \ell\nmid\operatorname{ord}(\psi),\qquad
 \psi_\ell^2=1,\qquad \psi(-1)=(-1)^e,\qquad
 \ell-1\nmid 2e,\\
 \ell\nmid |\operatorname{disc}K|\,h_K,\qquad
 \gcd(|\operatorname{disc}K|,N)=1,\\
 \ker\overline\chi_\ell\cap\ker\psi'\ \not\subseteq G_K .
\end{gathered}
\]
Here \(\overline\chi_\ell\) is the mod-\(\ell\) cyclotomic character,
\(h_K\) is the class number, and the last line uses Galois
characters on \(G_{\mathbb Q}\). The prime-to-\(\ell\) order assumption
makes \(\psi'\) precisely the character called \(\mathrm{chiP}\)
in the source. In particular \(\ell\ge5\).

Put \(\Gamma=G_{\mathbb Q}\), let \(\eta\) be the quadratic character
with kernel \(G_K\), and put \(\Psi=\chi_\ell^e\psi\),
\(\alpha=\Psi^{-1}\), and \(W=(F/\mathcal O)(\alpha)\).
Let \(\mathcal D\) consist of the decomposition groups at places above
primes dividing \(\ell N\) and the inertia groups at all finite primes.
Define
\[
 V=\ker\left(H^1_{\rm cont}(\Gamma,W)\longrightarrow
                      \prod_{D\in\mathcal D}H^1_{\rm cont}(D,W)\right).
\]
Fix a field isomorphism \(\iota:\overline{\mathbb Q}_\ell\to\mathbb C\)
and an element \(\sigma\in G_K\). Set
\[
 t=\Psi(\sigma)+1+\Psi(\sigma)^{-1},\qquad
 v_\ell(z)=-\log_\ell |z|_\ell\quad(z\ne0).
\]
The admissible \(\rho:G_K\to{\rm GL}_3(\overline{\mathbb Q}_\ell)\)
are the two classes used in the integer estimate. The first consists
of irreducible representations of level \(N\), with the source's
weight tuple \((e,0,-e)\) and its unlinked local condition. Their
twists \(\rho\otimes\chi_\ell^{-1}\) satisfy the conjugate-self-dual
automorphy condition. In the second class,
\(\rho\simeq{\bf1}\oplus W_2|_{G_K}\), where \(W_2\) is an irreducible
two-dimensional \(\overline{\mathbb Q}_\ell\)-representation of
\(G_{\mathbb Q}\). The source requires its level condition
at \(N|\operatorname{disc}K|\), weights \((e,-e)\), determinant of
cyclotomic weight zero, a newform of weight \(2e+1\) attached to
\(W_2\otimes\chi_\ell^{-e}\), and the rational automorphic witness
called \(\mathrm{WitAQ}\). These are the predicates
\(\mathrm{Adm3U}_s\) and \(\mathrm{Adm2A}\), respectively. In
particular the second class still supplies three-dimensional traces.

Suppose \(\operatorname{tr}\rho(\sigma)\ne t\) for every admissible
\(\rho\). Under the deformation, rigidity and automorphy inputs specified
below, the source gives a finite set \(X\) of pairwise nonisomorphic
admissible representations such that \(V\) has finite length and
\begin{equation}
 \frac{\operatorname{length}_{\mathcal O}V}{e_\mathcal O}
 \ \le\ 2\sum_{\rho\in X}v_\ell(\operatorname{tr}\rho(\sigma)-t).
 \label{eq:single-def-ineq}
\end{equation}
In the prime sum, \(\sigma\) is specialized to the inverse of the
chosen Frobenius element. It is arbitrary in this assertion.

\subsection{The strict cocycle in a marked deformation}

Write \(c\) for a complex conjugation. Then \(c^2=1\) and
\(\Psi(c)=1\). For now take the marked datum supplied over an
unramified coefficient ring. It gives \(g_1\in\Gamma\) with
\(\Psi(g_1)^2-1\) a unit. At the selected group \(D_\ell\) over
\(\ell\) it gives \(s\in D_\ell\) with \(\Psi(s)^2-1\) a unit.
If \(u^2-1\) is a unit, then so are \(u-1\) and
\(u^{-1}-1=-u^{-1}(u-1)\). These two markings will be used for
different purposes.

Every class \(x\in H^1(\Gamma,W)\) has a unique cocycle \(f_x\)
representing it with \(f_x(g_1)=0\). In fact if \(z\) represents \(x\),
put
\[
 m=(\alpha(g_1)-1)^{-1}z(g_1),\qquad
 f_x(g)=z(g)-(\alpha(g)-1)m .
\]
Changing \(z\) by the coboundary of \(u\) changes \(m\) by \(u\).
A coboundary vanishing at \(g_1\) has zero parameter. Hence this
normalization is \(\mathcal O\)-linear and injects cohomology into
the locally constant functions on \(\Gamma\).

Let \(b\in\mathcal O\setminus\{0\}\), \(Q_b=\mathcal O/(b)\).
For \(x\in V[b]\) the function \(bf_x=f_{bx}\) vanishes identically.
The map \(y\bmod b\mapsto y/b\bmod\mathcal O\) identifies \(Q_b\)
with \((F/\mathcal O)[b]\). In these coordinates write the cocycle
as \(u_x:\Gamma\to Q_b\). On any \(D\in\mathcal D\) there exists
\(m_D\in F/\mathcal O\) with
\(f_x(g)=(\alpha(g)-1)m_D\) for \(g\in D\). In particular
\((\alpha(s)-1)bm_{D_\ell}=0\) and the factor in parentheses is a
unit. Thus \(bm_{D_\ell}=0\) and the local coboundary at \(\ell\) has a
parameter in \(Q_b\) itself. At elements \(g\) of the imposed
inertia kernel, the marked datum has \(\Psi(g)=1\) and \(g\) belongs
to a group in \(\mathcal D\). The same formula gives \(f_x(g)=0\).

Also \(0=f_x(c^2)=2 f_x(c)\), hence \(f_x(c)=0\).
The identities \(f_x(1)=0\) and
\(f_x(gh)=f_x(g)+\alpha(g)f_x(h)\) show first
\(f_x(g_1^{-1})=0\) and then \(f_x(cg_1^{-1}c)=0\).
These vanishings impose the two marked off-diagonal conditions in
the normalized representation.

Let \(J_3\) be the antidiagonal identity and use on \({\rm GL}_3\)
the involution \(M\mapsto J_3 M^{-{\rm t}} J_3\).
The polarized reducible representation of \(\Gamma\) has matrix part
\(\operatorname{diag}(\Psi,1,\Psi^{-1})\) and component \(\eta\).
Put
\[
 B(u)=\begin{pmatrix}0&0&0\\u&0&0\\0&-u&0\end{pmatrix}.
\]
The differential of the involution fixes \(B(u)\), since
\(-J_3 B(u)^{\rm t}J_3=B(u)\). Diagonal conjugation sends it to
\(B(\alpha(g)u)\). Thus \(\gamma_x(g)=B(u_x(g))\) is an adjoint
cocycle of weight \(-1\). Its \((2,1)\) entry recovers \(u_x\).
It vanishes on the prescribed kernel, satisfies the marked
conditions, and is a coboundary at \(D_\ell\). This gives an
injective additive map from \(V[b]\) to the normalized matrix
cocycles with coefficient module \(Q_b\).

Let \(R\) be the ring representing the marked global deformation
condition, \(r\) its universal representation and \(x_0:R\to\mathcal O\)
the reducible point. The ring \(\mathcal O\oplus Q_b\) has product
\((a,u)(a',u')=(aa',au'+a'u)\). Multiplying the base representation
by \(1+\gamma_x(g)\) gives a representation over this ring because
the cross-products of the infinitesimals vanish and the remaining
identity is exactly the adjoint cocycle equation. The local coboundary
can be lifted to a local change of basis by lifting its scalar
parameter to \(\mathcal O\). The local deformation condition includes
these conjugates of the base lift. Consequently the representing
property gives a point of the form \((x_0,D_x)\). Multiplicativity
gives the Leibniz rule
\(D_x(uv)=x_0(u)D_x(v)+x_0(v)D_x(u)\).
On each universal matrix entry the value is the corresponding
entry of \(\gamma_x(g)\operatorname{diag}(\Psi(g),1,\Psi(g)^{-1})\).
In particular these values determine \(u_x\).

\subsection{The saturated cotangent lattice}

Write \(\pi:A\to R\) for the map from the graded model, \(x_A=x_0\pi\),
\(I=\ker x_A\), and
\[
 I^{(2)}_{\rm sat}
   =\{a\in A:\ \exists\,0\ne d\in\mathcal O,\ da\in I^2\}.
\]
The weight \(-1\) part of the kernel maps to a lattice \(L_-\) in
\(A/I^{(2)}_{\rm sat}\). In the presentation let
\(F_1,\ldots,F_{r+1}\) be the equations of weight \(-1\), and put
\(C_-=L_-/\sum_i\mathcal O[F_i]\). Here \(r\) is the rank of
\(L_-\), and brackets indicate images in the saturated quotient.

We use the following particular lifting property of the graded
presentation. For the \(D_x\) above, their pullbacks to \(A\) are
reductions modulo \(b\) of \(\mathcal O\)-valued derivations at
\(x_A\) supported in weight \(-1\). This is the assertion constructed
in the source's \(\mathrm{ModelLower.main\_der}\). Its local argument
lifts the weight-\(-1\) coboundary parameter at \(\ell\), hence gives
a compatible integral local derivation. The presentation theorem then
extends that derivation across the weighted ambient variables. It
only prescribes their reductions on the global quotient. The
extension is a derivation on \(A\), whereas \(D_x\) on \(R\) can
have intrinsically torsion coefficients.

For an integral lift \(\widetilde D_x\), Leibniz gives
\(\widetilde D_x(I^2)=0\). If \(df\in I^2\) with \(d\ne0\),
then \(d\widetilde D_x(f)=0\) in the domain \(\mathcal O\), so
\(\widetilde D_x(f)=0\). The reduction kills \(I^{(2)}_{\rm sat}\).
Each \(F_i\) maps to zero in \(R\), so \(D_x\pi(F_i)=0\).
The resulting map
\[
                 V[b]\longrightarrow
            \operatorname{Hom}_{\mathcal O}(C_-,Q_b)
\]
is injective. Indeed agreement on \(C_-\) gives agreement on all
of \(A_{-1}\). The pullbacks vanish on the other homogeneous pieces
by the support of their integral lifts, so they agree on \(A\).
Every universal matrix coordinate has a homogeneous lift to \(A\).
The formula for the \((2,1)\) coordinate now recovers \(u_x\), hence
\(x\). In particular saturation here uses torsion-free target
coefficients before reduction.

The presentation's spanning condition says that the classes of the
\(F_i\) span \(L_-\otimes_{\mathcal O}F\), and \(L_-\) is finite free.
Over the DVR write \(C_-\simeq\bigoplus_i\mathcal O/\varpi^{c_i}\).
If \(v_{\mathcal O}(b)=j\), its Hom group above has length
\(\sum_i\min(c_i,j)\), since the image of each cyclic generator can
be any element of \(Q_b\) killed by \(\varpi^{c_i}\).
An injection of finite sets of modules of finite length over
\(\mathcal O\) compares their lengths, since each module of length \(k\)
has \((\#\mathcal O/\varpi)^k\) elements by a composition series.
Thus \(\operatorname{length}V[b]\le\sum_i c_i\).
Every continuous cocycle into the discrete group \(W\) has finite
image because \(\Gamma\) is compact. Some power of \(\varpi\) kills
that image. Thus \(V=\bigcup_j V[\varpi^j]\). The lengths in this
increasing chain are bounded integers, and equality of two successive
finite lengths implies equality of the corresponding nested modules.
After its last strict increase the chain is constant. We conclude
\begin{equation}
                   \operatorname{length}_{\mathcal O} V
                        \le \sum_i c_i .
 \label{eq:single-cot-lower}
\end{equation}

\subsection{The two multiplicities in the intersection estimate}

We spell out the source of the coefficient \(2\) in
\eqref{eq:single-def-ineq}. Lift the three diagonal coordinates of
\(r(\sigma)\) to elements of \(A_0\) and let \(a\in A_0\) be their
sum minus \(t\), which belongs to \(\mathcal O\) in our normalization.
Then \(x_A(a)=0\). At each point \(f\) of the deformation ring,
\[
                 f(\pi(a))=\operatorname{tr}\rho_f(\sigma)-t
\]
because \(\eta(\sigma)=1\). Scaling a homogeneous coordinate of
weight \(k\) by \(z^k\), \(z\in\overline{\mathbb Q}_\ell^\times\),
fixes this trace. Call a geometric point mixed if both a positive
and a negative weight have some coordinate with nonzero value.
Take \(T\) meeting every mixed orbit on the equations. The
point and rigidity results mentioned below yield a finite such set.

One may change generators among the \(r+1\) weight \(-1\) equations.
Smith normal form gives generators with unchanged span and elements
\(y_i\in A_{-1}\) such that
\[
  F_i\equiv\varpi^{c_i}y_i\pmod{I^{(2)}_{\rm sat}}
     \quad(1\le i\le r),\qquad
  F_{r+1}\equiv0\pmod{I^{(2)}_{\rm sat}} .
\]
The other equations come in \(s\) of weight \(-2\), \(r'\) of weight
\(1\), \(s'\) of weight \(2\), and \(m\) of weight zero. Denote them
by \(G_i,E_i,E'_i,Z_i\), respectively. The chart used for the
intersection kills the \(Z_i\). One can replace \(a\) by a
weight-zero representative modulo the equations whose image in the
chart is nonzero. Its value and hence its valuation at every point
of \(T\) is unchanged.

For clarity the graded ring in the comparison is the three-variable
Rees algebra \(\mathcal R\) over \(A_0\) of the following
\(A_0\)-submodules of \(A\),
\[
    A_{-2},\quad A_2,\quad A_0F_{r+1}^2+A_{-2}^2 A_2.
\]
Its component at \(d=(d_0,d_1,d_2)\in\mathbb N^3\) has coefficients
in the product of the respective \(d_i\)-th powers. Let \(s_*\)
denote the following sequence of homogeneous sections. In degree
\((1,0,0)\) take \(F_1^2,\ldots,F_r^2\) and all \(G_i\). In
\((0,1,0)\) take all \(E_i^2,E'_i\). In \((0,0,1)\) take
\(F_{r+1}^2\). Let \(\mathcal J\) consist of polynomials whose
coefficients are killed by the chart map. It is a homogeneous
ideal. Write \(M=\mathcal R/\mathcal J\) and \(\chi(-;-)_d\)
for the Euler characteristic of the prescribed homogeneous sections
on the indicated module in multidegree \(d\), using lengths.

The lower chart estimate and upper point estimate in the source
give the following chain for all sufficiently large multidegrees
in the product order, with \(v_{\mathcal O}=e_{\mathcal O}v_\ell\):
\begin{align}
  2^{r+r'}\sum_i c_i
       &\le \chi(s_*;M/aM)_d
        \le \chi(a;M/(s_*)M)_d \notag\\
       &\le 2^{r+1+r'}\sum_{\xi\in T}v_{\mathcal O}(\xi(a)).
       \label{eq:single-rees}
\end{align}
The middle comparison uses regularity of \(a\) on \(M\). Indeed, in
each coefficient the chart sends it to a nonzero power series in a
domain, so \(a h\in\mathcal J\) implies \(h\in\mathcal J\). Applying
the nonzerodivisor rule for Euler characteristics, permuting the
sections, and applying the upper bound by the quotient by \(s_*\)
gives the stated comparison. The source proves the required
admissibility for this sequence before it uses these Euler rules.

Here is the coefficient calculation on the lower side. The leading
chart quotient modulo \(\varpi\) has as Hilbert function the number
of monomials with total negative weight \(2d_0+4d_2\) and total
positive weight \(2d_1+2d_2\). This identification and the lower
Euler bound on pulling out scalar factors from homogeneous sections
are inputs of the chart argument. In this quotient the first \(r\)
sections have leading parts \(\varpi^{2c_i}y_i^2\). The latter Euler
bound gives \(\sum_i2c_i\) times a finite difference of the Hilbert
function. The difference uses \(r-1+s\) steps in direction zero,
\(r'+s'\) in direction one, and one in direction two.

We compute this difference directly. The chart after cutting the
weight-zero equations has \(r,s,r'+1,s'+1\) variables of weights
\(-1,-2,1,2\), respectively. Thus its bivariate monomial generating
function is
\[
 H(X,Y)=\frac{1}{(1-X)^r(1-X^2)^s(1-Y)^{r'+1}(1-Y^2)^{s'+1}} .
\]
The coefficient is read at \(X^{2d_0+4d_2}Y^{2d_1+2d_2}\).
The differences therefore multiply \(H\) by
\((1-X^2)^{r-1+s}(1-Y^2)^{r'+s'}(1-X^4Y^2)\).
Writing \(P=(1+X)^{r-1}(1+Y)^{r'}\), the product equals
\[
 \frac{P}{(1-X)(1-Y)}
   +\frac{PY^2(1+X+X^2+X^3)}{(1-Y)(1-Y^2)} .
\]
This identity follows by splitting \(1-X^4Y^2=(1-Y^2)+Y^2(1-X^4)\).
The second summand has bounded \(X\)-degree. At large degrees in
both variables the first summand has coefficient
\(P(1,1)=2^{r-1+r'}\). Multiplying this by \(\sum_i2c_i\)
gives exactly the leftmost side of \eqref{eq:single-rees}.

On the upper side the generic local length of each weight piece at
the restriction of a mixed point to \(A_0\) is at most one for the
original equations. This uses the point rigidity and the graded
slice theorem. Replacing a homogeneous equation \(u\) of weight
\(k\) by its square at most doubles a uniform bound on the local
lengths of all the weight pieces. Indeed in a graded quotient \(P'\)
the map \(P'/u^2P'\to P'/uP'\) has kernel \(uP'/u^2P'\). The latter
is a quotient of \(P'/uP'\) via multiplication by \(u\), with a
shift by \(k\). This remains a right-exact sequence on each graded
piece and on localization. Thus a weight-\(j\) piece of the squared
quotient has length at most the sum for weights \(j\) and \(j-k\)
of the original quotient. Iterating over the \(r+1\) equations \(F_i\)
and the \(r'\) equations \(E_i\) gives \(2^{r+1+r'}\). The other
equations are unchanged. The upper Euler estimate uses a surjection
onto the relevant large Rees pieces, then the length-versus-valuation
formula on their support. Only the restrictions of mixed points can
contribute. Their values of \(a\) are nonzero and have nonnegative
valuation. Summing over \(T\) majorizes the sum over distinct
restrictions even when several elements restrict to the same point.

The algebraic results used for these two estimates occur in
\(\mathrm{AdmissibleMonoIdeal.low\_holds}\),
\(\mathrm{TwoSigned.up\_holds}\), and
\(\mathrm{AlgMemCoef.core\_holds}\). The calculations above prove
the two powers of two and the regularity check. They use, as separate
inputs, the chart identification, the Euler rules and support
comparison, and the generic local-length-one assertion. Since
\(r\ge1\), cancellation of the positive factor \(2^{r+r'}\) in
\eqref{eq:single-rees} yields
\[
               \sum_i c_i\le
                 2\sum_{\xi\in T}v_{\mathcal O}(\xi(a)).
\]

\subsection{Identification of the points and coefficient descent}

The cohomological hypothesis used for the construction consists of
the local deformation inputs, normalized Poitou--Tate duality, global
cohomological finiteness and the global Euler characteristic theorem,
and the size bound for \(H^2\) called Lemma~\(K\).
Together with the presentation theorem called Lemma~\(P\), these
give the marked ring and its graded model, including the integral
lifting property used above. Over unramified coefficients the
automorphy input supplies the marking, finiteness of the subalgebra
of weight-zero matrix monomials over \(\mathcal O\), automorphy for
points of the first type, and modularity for the two-dimensional
extensions occurring at points of the second type.

More precisely, the point condition is deduced in
\(\mathrm{FPtHypU}\). It identifies a point that is both upper and
lower with the base point, makes the prescribed tangent vectors at
the base coboundaries by nonpositive-weight Bloch--Kato vanishing,
and makes tangent vectors at mixed points coboundaries by adjoint
Selmer vanishing. For points of the first type the last vanishing
uses \(\mathrm{RigidityFirstCase}\) at the automorphic representation.
For points of the second type, the proof uses the rank-two Kato input
\(\mathrm{KatoTwoW}\), \(\mathrm{RigiditySecondCase}\) for the
two-dimensional extension of determinant \(\eta\), and its modularity.
These give the point, tangent and spanning hypotheses of the graded
intersection estimate. The first type is admissible by
\(\mathrm{AdmI}_s\). The second is admissible by the theorem on
modular extensions of type II. These are arithmetic inputs, additional
to the linear and intersection calculations above.

The local-point comparison in the graded presentation puts every
mixed point, up to the scaling action, over a point of \(R\). Conversely
it identifies mixed points of \(R\) with those graded points having
both signs. The matrix coordinates of weight \(j\) scale by \(z^j\). On representations this is conjugation by
\(\operatorname{diag}(z,1,z^{-1})\). Choose one point of \(R\) per
mixed orbit. The marked shape theorem says that for two mixed points
an isomorphism between their restrictions to \(G_K\) is induced by
such a diagonal conjugacy. Thus restriction to \(G_K\) gives a set
\(X\) with one representative per isomorphism class in its image,
and summing over these representatives preserves the valuation
sum. Each term is the trace difference above. Combining the upper
estimate with \eqref{eq:single-cot-lower} and dividing by
\(e_\mathcal O\) proves \eqref{eq:single-def-ineq} for the
unramified coefficient choice.

Finally the order \(m_\psi\) of \(\psi\) is prime to \(\ell\).
Choose \(f>0\) with \(m_\psi\mid\ell^f-1\). The required roots
of unity lift by Hensel's lemma from \(\mathbb F_{\ell^f}\) to
the unramified degree-\(f\) extension, so the same character values
occur in an unramified coefficient ring. To compare with the original
ring take a finite compositum. For an inclusion of coefficient rings
\(\mathcal O_1\subset\mathcal O_2\) the latter is finite free, and
\((F_1/\mathcal O_1)\otimes_{\mathcal O_1}\mathcal O_2
 \simeq F_2/\mathcal O_2\).
After choosing an \(\mathcal O_1\)-basis of \(\mathcal O_2\),
continuous cochains and restrictions for this scalar extension are
finite direct sums of their originals, so the kernels defining
the strict group commute with extension. For a finite length
module a composition series shows that its length is multiplied by
the relative ramification index, since the extension of a residue
field factor is \(\mathcal O_2/\varpi_1\). Finiteness also descends
by finite freeness. The denominator \(v_{\mathcal O_i}(\ell)\)
changes by the same index. Thus the normalized length, its finiteness,
and the target trace are the same in both coefficient realizations.
This is the coefficient comparison of \(\mathrm{CoefIndep}\) in the
proof of \(\mathrm{FOneCharU}\).
